\begin{theorem}[C.2.1 (Known, with derivation): Lead 6's formula]
Let $M=M_\phi$ have the orientation above. Let $W_8$ be a smooth
compact oriented spin$^c$ eight-manifold with $\partial W_8=M$.
Let $c\in H^2(W_8;\Z)$ be the first Chern class of its determinant
line, and choose a trivialization on the boundary which induces the
spin structure used on $M$. Write $p=p_1(TW_8)$, choose a relative
lift $\bar c$, and let $\bar p$ be the unique rational relative lift
of $p$. Then, with all products evaluated on $[W_8,M]$,
\begin{equation}\label{eq:c-spin-c}
 \mu(M)\equiv
 \frac{\bar p^{,2}-4\sigma(W_8)}{896}
 -\frac{\bar p\bar c^{,2}}{192}
 +\frac{\bar c^{,4}}{384}\pmod{\Z}.
\end{equation}
The value is independent of the lift $\bar c$. If $W_8$ is almost
complex, use its canonical spin$^c$ structure, for which
$c=c_1(TW_8)$ and $p=c_1^2-2c_2$. If $W_8$ is spin and its
determinant line is trivial, this reduces to
\[
 28\mu(M)=\frac{p_1^2-4\sigma(W_8)}{32}\pmod{28}.
\]
All the boundary cohomology and spin-structure conditions in this
statement hold for the mapping torus; the existence of the particular
geometric $W_8$ proposed in Part B is a separate assertion.
\end{theorem}

\begin{proof}
The mapping-torus exact sequence, or Proposition C.1.1, gives
$H^k(M;\Z)=0$ for $2\leq k\leq5$ and $H^1(M;\Z)=\Z$.
The same vanishing in degree two holds with $\Z/2$ coefficients,
so $M$ is spin. Its spin structures form a torsor for
$H^1(M;\Z/2)=\Z/2$. The exact sequence of the pair gives an
isomorphism
$H^4(W_8,M;\Q)\to H^4(W_8;\Q)$ and a surjection
$H^2(W_8,M;\Z)\to H^2(W_8;\Z)$, proving the assertions
about lifts and determinant trivializations.

If two lifts of $c$ differ by $\delta\alpha$ with
$\alpha\in H^1(M;\Z)$, the relative cup-product identity gives
$\delta\alpha\smile b=\delta(\alpha\smile i^*b)$.
The restrictions of $c$ and $p$ vanish on $M$. Every term in the
change of $\bar p\bar c^2$ or $\bar c^4$ is therefore zero;
in particular $(\delta\alpha)^2=0$. This proves lift independence.
There is no assertion here that the relative class $\bar c$ itself
is unique.

For a closed spin$^c$ eight-manifold the degree-eight part of the
Dirac index density is
\[
 [e^{c/2}\widehat A]_8
 =\frac{7p_1^2-4p_2}{5760}
       -\frac{p_1c^2}{192}+\frac{c^4}{384}.
\]
Eliminate $p_2$ using $\sigma=(7p_2-p_1^2)/45$.
The result is exactly the right side of
\eqref{eq:c-spin-c}. To justify its use on a boundary, glue $W_8$
to the negative of a spin coboundary inducing the same boundary spin
structure; extend the determinant line by its chosen trivialization.
The glued manifold is spin$^c$, so its Dirac index is an integer.
Novikov additivity and the relative-product comparison identify that
integer with the difference between \eqref{eq:c-spin-c} and the
spin formula for $M$. The spin formula is
\citet[Section 6, equation (11)]{EellsKuiper1962}.
The spin$^c$ version is also explicitly recorded, with normalization
$28\mu$ and characteristic class
$\widehat p=(p_1+c^2)/2$, by
\citet[Section 2.2, equation (2.1)]{CrowleyNordstrom2020}.
Substituting this $\widehat p$ gives the same coefficients; both
polynomial identities are verified exactly in \path{code/c_exact.py}.

A spin coboundary can be supplied using the connected-sum description:
bound the product by $S^1\times D^7$ and bound the sphere summand
by its spin coboundary. Thus no inference from simple connectivity of
$M$ is made; indeed $\pi_1(M)=\Z$. The original Eells--Kuiper
paper's unstructured formulation has an extra condition on
$H^1(W;\Z/2)\to H^1(M;\Z/2)$ to ensure spin structures glue.
Our explicit requirement that the boundary spin structures match
supplies precisely that gluing datum. By Proposition C.1.1 the two
possible spin structures on this particular $M$ have equal invariant.
Finally the splitting principle gives
$p_1(TW_8)=c_1^2-2c_2$ for an almost complex tangent bundle.
\end{proof}

\subsection{C.2 supplement. A computed coboundary with a different boundary}

